\documentclass[../../main.tex]{subfiles}

\begin{document}

\section{习题}\label{section:ch10-习题}

\begin{exercise}\label{exercise:ch10-对偶粘性解}
假设 $u$ 是 $u_t+H(Du,x)=0$ 在 $\mathbb{R}^n\times(0,\infty)$ 中的粘性解。证明 $\widehat{u}\coloneqq-u$ 是
\begin{align*}
\widehat{u}_t+\widehat{H}(D\widehat{u},x)=0 \quad \text{在 } \mathbb{R}^n\times(0,\infty) \text{ 中}
\end{align*}
的粘性解，其中 $\widehat{H}(p,x)\coloneqq-H(-p,x)$。
\end{exercise}

\begin{exercise}\label{exercise:ch10-一致极限}
设 $\{u_k\}_{k=1}^\infty$ 是Hamilton--Jacobi方程
\begin{align*}
(u_k)_t+H(Du_k,x)=0 \quad \text{在 } \mathbb{R}^n\times(0,\infty) \text{ 中}
\end{align*}
的粘性解（$k=1,2,\cdots$），假设 $u_k\to u$ 一致。还假设 $H$ 连续。证明 $u$ 是
\begin{align*}
u_t+H(Du,x)=0 \quad \text{在 } \mathbb{R}^n\times(0,\infty) \text{ 中}
\end{align*}
的粘性解。因此粘性解的一致极限仍是粘性解。
\end{exercise}

\begin{exercise}\label{exercise:ch10-抛物正则化无关}
假设对每个 $\varepsilon>0$，$u^\varepsilon$ 是抛物型方程
\begin{align*}
u_t^\varepsilon+H(Du^\varepsilon,x)-\varepsilon\sum_{i,j=1}^n a_{ij}u_{x_i x_j}^\varepsilon=0 \quad \text{在 } \mathbb{R}^n\times(0,\infty) \text{ 中}
\end{align*}
的光滑解，其中光滑系数 $a_{ij}$（$i,j=1,\cdots,n$）满足\cref{chapter:ch06-二阶椭圆型方程}中的一致椭圆性条件。还假设 $H$ 连续，且 $u^\varepsilon\to u$ 一致当 $\varepsilon\to 0$。证明 $u$ 是 $u_t+H(Du,x)=0$ 的粘性解。（本习题表明粘性解不依赖于抛物正则化的精确结构。）
\end{exercise}

\begin{exercise}\label{exercise:ch10-压缩性质}
设 $u_i$（$i=1,2$）是
\begin{align*}
\begin{cases}
(u_i)_t+H(Du_i,x)=0 & \text{在 } \mathbb{R}^n\times(0,\infty) \text{ 中},\\
u_i=g_i & \text{在 } \mathbb{R}^n\times\{t=0\} \text{ 上}
\end{cases}
\end{align*}
的粘性解。假设 $H$ 满足 \S10.2 中的条件\eqref{align:ch10-H的Lipschitz条件}。证明压缩性质
\begin{align*}
\sup_{\mathbb{R}^n}|u_1(\cdot,t)-u_2(\cdot,t)|\leqslant\sup_{\mathbb{R}^n}|g_1-g_2| \quad (t\geqslant 0).
\end{align*}
\end{exercise}

\begin{exercise}\label{exercise:ch10-距离函数粘性解}
\begin{enumerate}[(a)]
\item 证明 $u(x)\coloneqq 1-|x|$ 是
\begin{align}
\begin{cases}
|u'|=1 & \text{在 } (-1,1) \text{ 中},\\
u(-1)=u(1)=0
\end{cases}
\label{align:ch06-正则PDE}
\end{align}
的粘性解。这意味着对每个 $v\in C^\infty(-1,1)$，若 $u-v$ 在点 $x_0\in(-1,1)$ 处有最大值（最小值），则 $|v'(x_0)|\leqslant 1$（$\geqslant 1$）。
\item 证明 $\widetilde{u}(x)\coloneqq |x|-1$ 不是\eqref{align:ch06-正则PDE} 的粘性解。
\item 证明 $\widetilde{u}$ 是
\begin{align}
\begin{cases}
-|\widetilde{u}'|=-1 & \text{在 } (-1,1) \text{ 中},\\
\widetilde{u}(-1)=\widetilde{u}(1)=0
\end{cases}
\label{align:ch10-PDE星星}
\end{align}
的粘性解。（提示：\eqref{align:ch10-PDE星星} 的粘性解是什么意思？）
\item 为什么\eqref{align:ch06-正则PDE}、\eqref{align:ch10-PDE星星} 有不同的粘性解？
\end{enumerate}
\end{exercise}

\begin{exercise}\label{exercise:ch10-eikonal方程}
设 $U\subset\mathbb{R}^n$ 开、有界。令 $u(x)\coloneqq\operatorname{dist}(x,\partial U)$（$x\in U$）。证明 $u$ Lipschitz 连续，且它是 eikonal 方程
\begin{align*}
|Du|=1 \quad \text{在 } U \text{ 中}
\end{align*}
的粘性解。这意味着对每个 $v\in C^\infty(U)$，若 $u-v$ 在点 $x_0\in U$ 处有最大值（最小值），则 $|Dv(x_0)|\leqslant 1$（$\geqslant 1$）。
\end{exercise}

\begin{exercise}\label{exercise:ch10-间断面条件}
假设开集 $U\subset\mathbb{R}^n$ 被光滑超曲面 $\Gamma$ 分割为子区域 $V^+$ 和 $V^-$。设 $\nu$ 表示 $\Gamma$ 的单位法向量，指向 $V^+$ 内部。假设 $u$ 是
\begin{align*}
H(Du)=0 \quad \text{在 } U \text{ 中}
\end{align*}
的粘性解，且 $u$ 在 $\overline{V^+}$ 和 $\overline{V^-}$ 中光滑。记 $u_\nu^+$ 为沿 $\Gamma$ 从 $V^+$ 内部取的 $Du\cdot\nu$ 的极限，记 $u_\nu^-$ 为从 $V^-$ 内部取的极限。证明沿 $\Gamma$ 我们有不等式
\begin{align*}
H(\lambda Du^-+(1-\lambda)Du^+)\geqslant 0 \quad \text{若 } u_\nu^-\leqslant u_\nu^+,
\end{align*}
且
\begin{align*}
H(\lambda Du^-+(1-\lambda)Du^+)\leqslant 0 \quad \text{若 } u_\nu^+\leqslant u_\nu^-,
\end{align*}
对每个 $0\leqslant\lambda\leqslant 1$，其中 $Du^\pm$ 表示沿 $\Gamma$ 从 $V^\pm$ 取的梯度。
\end{exercise}

\begin{exercise}\label{exercise:ch10-从明暗恢复形状}
由 $u:\mathbb{R}^2\to\mathbb{R}$ 的图像描述的曲面被来自竖直方向 $e_3$ 的平行光线照射。我们假设曲面具有常数反照率，此外是 Lambertian 的，即入射光线在所有方向上均匀散射。则在点 $x\in\mathbb{R}^2$ 上方反射光的强度 $i=i(x)$ 由公式 $i=e_3\cdot\nu$ 给出，其中 $\nu$ 是曲面指向上方的单位法向量。证明 $u$ 求解形式为
\begin{align*}
|Du|=n
\end{align*}
的偏微分方程，其中 $n=n(x)$ 是给定函数，由强度 $i$ 计算得到。（通过求解此偏微分方程求 $u$ 来确定曲面，是\emph{从明暗恢复形状}问题。）
\end{exercise}

\begin{exercise}\label{exercise:ch10-追及问题}
一艘游艇从正 $x_1$ 轴上的点 $(x_1,0)$ 出发，以速度 $b_1>0$ 向右航行。另一艘游艇最初在正 $x_2$ 轴上的点 $(0,x_2)$ 处，开始追及，始终以速度 $b_2>b_1$ 朝向第一艘游艇航行。求
\begin{align*}
u(x_1,x_2)\coloneqq \text{第二艘游艇拦截第一艘游艇所需的时间}
\end{align*}
所求解的偏微分方程。（将此视为一个动态规划问题，但没有控制。）
\end{exercise}

\begin{exercise}[无限时域控制问题。]\label{exercise:ch10-无限时域控制}假设 $\mathbf{f}$ 和 $r$ 满足 \S10.3 中给出的条件。给定点 $x\in\mathbb{R}^n$ 和属于 $\mathcal{A}\coloneqq\{\boldsymbol{\alpha}:[0,\infty)\to A\mid\boldsymbol{\alpha}(\cdot)\text{ 可测}\}$ 的控制，令 $\mathbf{x}(\cdot)$ 为常微分方程
\begin{align*}
\begin{cases}
\dot{\mathbf{x}}(s)=\mathbf{f}(\mathbf{x}(s),\boldsymbol{\alpha}(s)) & (s>0),\\
\mathbf{x}(0)=x
\end{cases}
\end{align*}
的唯一解。固定 $\lambda>0$，定义折扣代价
\begin{align*}
C_x[\boldsymbol{\alpha}(\cdot)]\coloneqq\int_0^\infty e^{-\lambda s}r(\mathbf{x}(s),\boldsymbol{\alpha}(s))\,\mathrm{d}s.
\end{align*}
定义值函数
\begin{align*}
u(x)\coloneqq\inf_{\boldsymbol{\alpha}(\cdot)\in\mathcal{A}}C_x[\boldsymbol{\alpha}(\cdot)].
\end{align*}
\begin{enumerate}[(a)]
\item 证明 $u$ 有界，且若 $\lambda>\operatorname{Lip}[\mathbf{f}]$，则 $u$ Lipschitz 连续。
\item 证明若 $0<\lambda\leqslant\operatorname{Lip}[\mathbf{f}]$，则 $u$ 对某个指数 $0<\alpha<1$ 是 H\"older 连续的。
\end{enumerate}
\end{exercise}

\begin{exercise}[续。]\label{exercise:ch10-无限时域HJB}证明值函数 $u$ 是偏微分方程
\begin{align*}
\lambda u-\min_{a\in A}\{\mathbf{f}(x,a)\cdot Du+r(x,a)\}=0 \quad \text{在 } \mathbb{R}^n \text{ 中}
\end{align*}
的粘性解。（这意味着若 $v$ 光滑且 $u-v$ 在点 $x_0$ 处有局部最大值，则
\begin{align*}
\lambda u-\min_{a\in A}\{\mathbf{f}(x,a)\cdot Dv+r(x,a)\}\leqslant 0
\end{align*}
在 $x_0$ 处，且若 $u-v$ 在 $x_0$ 处有局部最小值，则相反的不等式成立。）

接下来的一系列习题发展了二阶完全非线性椭圆型偏微分方程粘性解的一些理论。
\end{exercise}

\begin{exercise}\label{exercise:ch10-椭圆性定义}
回忆 \S A.1，若 $R,S\in\mathbb{S}^n$，我们写 $R\geqslant S$ 若 $R-S$ 半正定。函数 $F:\mathbb{S}^n\times\mathbb{R}^n\times\mathbb{R}^n\to\mathbb{R}$，$F=F(R,p,x)$，称为\emph{椭圆的}，只要
\begin{align*}
R\geqslant S \text{ 蕴含 } F(R,p,x)\leqslant F(S,p,x).
\end{align*}
这里 $\mathbb{S}^n$ 表示实 $n\times n$ 对称矩阵空间。
\begin{enumerate}[(a)]
\item 证明 $F(R)=-\operatorname{tr}R$ 是椭圆的。
\item 更一般地，证明若 $A\in\mathbb{S}^n$ 且 $A\geqslant 0$，则 $F(R)=-A:R=-\operatorname{tr}(AR^T)$ 是椭圆的。
\item 证明若对每个 $k=1,\cdots,m$，$F_k$ 是椭圆的，则 $\max_k F_k$ 和 $\min_k F_k$ 也是椭圆的。
\end{enumerate}
\end{exercise}

\begin{exercise}\label{exercise:ch10-完全非线性粘性解}
设 $F$ 连续且椭圆。我们称函数 $u\in C(U)$ 为完全非线性椭圆型偏微分方程
\begin{align}
F(D^2u,Du,x)=0 \quad \text{在 } U \text{ 中}
\label{align:ch10-完全非线性PDE}
\end{align}
的粘性解，只要对每个 $v\in C^\infty(U)$，(i) 若 $u-v$ 在点 $x_0\in U$ 处有局部最大值，则 $F(D^2v(x_0),Dv(x_0),x_0)\leqslant 0$；(ii) 若 $u-v$ 在点 $x_0\in U$ 处有局部最小值，则 $F(D^2v(x_0),Dv(x_0),x_0)\geqslant 0$。证明若 $u$ 是\eqref{align:ch10-完全非线性PDE} 的 $C^2$ 解，则 $u$ 是粘性解。
\end{exercise}

\begin{exercise}\label{exercise:ch10-椭圆一致极限}
假设 $u_k$ 是椭圆型偏微分方程
\begin{align*}
F(D^2u_k,Du_k,x)=0 \quad \text{在 } U \text{ 中}
\end{align*}
的粘性解，$k=1,2,\cdots$。假设 $u_k\to u$ 一致，证明 $u$ 是
\begin{align*}
F(D^2u,Du,x)=0 \quad \text{在 } U \text{ 中}
\end{align*}
的粘性解。
\end{exercise}

\end{document}